% ============================================================
% Chapter 4: Proposed Method: Specification and Design
%   Section: Data Architecture
% Included from chapters/04-proposed-method/proposed-method.tex
% Sources: docs/decisions/storage.md, docs/decisions/graph.md.
% ============================================================

\section{Data Architecture}
\label{sec:data-architecture}

Three stores hold what the system produces, each for a different use.

\paragraph{A three-layer data lake.} Every analysis is written into a
lake of three layers, one write per stage as that stage completes:

\begin{itemize}
    \item \textbf{Raw}: the extracted article and its fetch metadata.
    Nothing model-derived, so it stays valid when every model changes.
    \item \textbf{Processed}: the full enriched article and its
    fact-check report, rejections and discarded evidence included. It
    is the layer everything else is re-derived from.
    \item \textbf{Exploitation}: one flat, ready-to-serve record per
    analysis, with the editorial decision precomputed as
    \emph{publishable} --- admitted, and with no claim judged
    \texttt{FALSE} or \texttt{MISLEADING}.
\end{itemize}

Writing per stage means that a fact-check that fails (the search
engine down, the model timing out) does not discard an enrichment that
already cost the whole NLP stack. Every record carries its
\emph{lineage}: the run, the article, a hash of the body that joins the
layers, its parent record, the pipeline version, the code revision, the
model versions and any threshold the run overrode. Record identifiers
are derived from the layer, the article and the run, so repeating a
write overwrites rather than duplicates it, and an append-only manifest
keeps the history of every write.

\paragraph{A vector store.} Qdrant holds one embedding per analysed
article. It serves three purposes: duplicate detection at admission,
the system's own corpus searched for evidence, and the related-article
signal. An article is added to it only after its fact-check, so it
cannot be found as evidence for its own claims.

\paragraph{A graph.} Every finished analysis is also written to Neo4j,
in one transaction per article, after the lake. Its schema
(Figure~\ref{fig:graph-schema}) has seven node types: articles, the
sources that publish them and the evidence pages, the entities they
mention, their topics, their claims and the five verdicts. Every
relationship carries the method that wrote it --- the pipeline or a
manual label --- so a hand label and the model's verdict on the same
claim sit side by side instead of one overwriting the other. The lake
is the record and the graph a view of it: a graph write that fails
costs one event and never the analysis, and a synchronisation endpoint
rebuilds the graph from the lake and the labelled facts at any time.

\begin{figure}[ht]
\centering
\begin{adjustbox}{max width=\linewidth}
\begin{tikzpicture}[
    n/.style={draw, rounded corners=6pt, font=\footnotesize, fill=ganttplan!15,
              minimum width=20mm, minimum height=7mm},
    arr/.style={-{Stealth[length=2mm]}, thick},
    lbl/.style={font=\scriptsize\ttfamily, fill=white, inner sep=1.5pt},
]
\node[n] (top) at (0, 2.2)  {Topic};
\node[n] (art) at (0, 0)    {Article};
\node[n] (src) at (6, 0)    {Source};
\node[n] (ent) at (-6, 0)   {Entity};
\node[n] (cl)  at (0, -2.4) {Claim};
\node[n] (ver) at (-3.5, -4.6) {Verdict};
\node[n] (ev)  at (3.5, -4.6)  {Evidence};
\draw[arr] (art) -- node[lbl]{ABOUT} (top);
\draw[arr] (art) -- node[lbl]{PUBLISHED\_BY} (src);
\draw[arr] (art) -- node[lbl]{MENTIONS} (ent);
\draw[arr] (art) -- node[lbl]{CONTAINS\_CLAIM} (cl);
\draw[arr] (cl) -- node[lbl]{HAS\_VERDICT} (ver);
\draw[arr] (cl) -- node[lbl]{CHECKED\_AGAINST} (ev);
\draw[arr] (cl) -| node[lbl, pos=0.25]{MENTIONS} (ent);
\draw[arr] (ev) -| node[lbl, pos=0.25]{PUBLISHED\_BY} (src);
\end{tikzpicture}
\end{adjustbox}
\caption{The graph schema. Every relationship carries
\texttt{method} (\texttt{pipeline} or \texttt{manual}).}
\label{fig:graph-schema}
\end{figure}
